% $Header: /cvsroot/latex-beamer/latex-beamer/examples/beamerexample2.tex,v 1.7 2004/10/07 20:53:07 tantau Exp $

% This file is included by beamerexample2.article.tex and
% beamerexample2.beamer.tex 

% Copyright 2003 by Till Tantau $<$tantau@cs.tu-berlin.de$>$.
%
% This program can be redistributed and/or modified under the terms
% of the LaTeX Project Public License Distributed from CTAN
% archives in directory macros/latex/base/lppl.txt.

%
% The purpose of this example is to demonstrate the usage of the
% nameslide command
%

% \mode<article>
% {
%   \usepackage{fullpage}
%   \usepackage{pgf}
%   \usepackage{hyperref}
%   \setjobnamebeamerversion{presentacion.beamer}
% }
% 
% \mode<presentation>
% {
%   \usetheme{Marzo}
% 
%   \setbeamercovered{transparent}
% }

\documentclass{beamer}

\usepackage{beamerthemesplit}
\usepackage[latin1]{inputenc}
\usepackage[spanish]{babel}

\setbeamertemplate{blocks}[rounded]
\usetheme{Madrid}

\title{Presentaci�n del GUL}
\author{Jes�s Espino Garc�a}
\date{3 de Marzo de 2008}
\subject{Presentaci�n del GUL}

\institute[GUL UC3M]{
  Grupo de usuarios de Linux\\
  Universidad Carlos III de Madrid.\\
  \ \\
  \includegraphics[height=2cm]{gul}}

\begin{document}

  \frame{\titlepage}
  
  \begin{frame}
    \frametitle{�Qu� es el GUL?}

    \begin{block}{El GUL}
    \begin{itemize}
      \item Es una asociaci�n de la Universidad Carlos III de Madrid.
      \item Es el Grupo de Usuarios de Linux de la Universidad Carlos III de Madrid.
      \item Dedicada a promover el uso del software libre.
      \item Consta de mas de 300 socios de todo el mundo.
      \item Ha colaborado en la organizaci�n de congresos de primera linea en
            el �mbito del software como son Hispalinux y Javahispano.
    \end{itemize}
    \end{block}
  \end{frame}
  
  \begin{frame}
    \frametitle{\begin{center}Organizaci�n.\end{center}}
  \end{frame}
  
  \begin{frame}
    \frametitle{Grupos}
    \begin{block}{Tres clases de grupos}
    \begin{itemize}
      \item Junta: Grupo encargado de la toma de decisiones.
      \item Grupos t�cnicos: Grupos encargados de llevar a buen puerto los proyectos del GUL.
      \item Grupos de inter�s: Grupos con intereses comunes.
    \end{itemize}
    \end{block}
  \end{frame}

  \begin{frame}
    \frametitle{Grupos}
    \begin{block}{La Junta}
    \begin{itemize}
      \item Un presidente: Jes�s Espino Garc�a.
      \item Un vicepresidente/secretario: Cl�udio Jim�nez.
      \item Un tesorero: Samu�l Rodr�guez Sevilla.
      \item Un vocal: Rafael Bravo.
      \item Un coordinador de actividades: Jose Mar�a Chumo.
    \end{itemize}
    \end{block}
  \end{frame}
  
  \begin{frame}
    \frametitle{Grupos}
    \begin{block}{Grupos t�cnicos}
    \begin{itemize}
      \item gul-root: Encargado de la administraci�n de los servidores del GUL.
      \item gul-www: Encargado de la pagina web del GUL.
      \item gul-myc: Encargado de marketing y comunicaci�n externas.
      \item gul-cursos: Encargado de la organizaci�n de los cursos.
    \end{itemize}
    \end{block}
    \begin{block}{Grupos de interes}
    \begin{itemize}
      \item gul-doc: Documentaci�n y traducci�n.
      \item gul-sec: Seguridad.
      \item gul-dev: Desarrollo.
    \end{itemize}
    \end{block}
  \end{frame}
  
  \begin{frame}
    \frametitle{\begin{center}Proyectos.\end{center}}
  \end{frame}
  
  \begin{frame}
    \frametitle{Proyectos}
    \begin{block}{Proyectos del �ltimo a�o}
    \begin{itemize}
      \item Cursos de Noviembre.
      \item Cursos de Marzo.
      \item Reuniones t�cnicas.
      \item Hospedaje de Gentoo-es.
      \item Mirror de Debian.
      \item Documentaci�n interna.
      \item Servidor de Proyectos.
      \item Cabaret.
      \item Reuniones del GUL.
    \end{itemize}
    \end{block}
  \end{frame}
  
  \begin{frame}
    \frametitle{Proyectos}
    \begin{block}{Proyectos futuros}
    \begin{itemize}
      \item Cursos de Noviembre.
      \item Cursos de Marzo.
      \item Cursos en Getafe.
      \item Reuniones t�cnicas.
      \item Hospedaje de Gentoo-es.
      \item Documentaci�n interna.
      \item Alta de nuevos servicios.
      \item Mejora de las infraestructuras de la asociacion
      \item Reuniones del GUL.
      \item Mirrors.
      \item \dots
    \end{itemize}
    \end{block}
  \end{frame}
  
  \begin{frame}
    \frametitle{Reuniones del GUL}
    \begin{block}{Objetivos}
    \begin{itemize}
      \item Toma de decisiones.
      \item Aportaci�n de ideas y proyectos.
      \item Revisi�n del estado de los proyectos.
      \item Incorporaci�n de nuevos activistas al GUL.
      \item \dots
    \end{itemize}
    \end{block}
    Se suelen hacer el segundo viernes de cada mes, pero lo normal es que alguien de la junta avise en la lista de correo.
  \end{frame}
  
  \begin{frame}
    \frametitle{Dudas}
    \dots
  \end{frame}
  
  \begin{frame}
    \frametitle{\begin{center}Fin\end{center}}
    \begin{center}
      \includegraphics[height=3cm]{gul}
    \end{center}
  \end{frame}

\end{document}
